\section{Probability Theory}

\subsection{Fundamental Conclusions in Probability}

\begin{proposition}[Mills' ratio]
For every $t>0$,
\[
\frac{t}{t^2+1}\phi(t)
\leq
\overline{\Phi}(t)
\leq
\frac{\phi(t)}{t}.
\]
Consequently, the logarithm of the Gaussian upper tail has leading term
$-t^2/2$ as $t\to\infty$, where
\[
\phi(t)=\frac{1}{\sqrt{2\pi}}e^{-t^2/2},
\qquad
\overline{\Phi}(t)
=\int_t^\infty \phi(x)\,\mathrm{d}x
=1-\Phi(t).
\]
\end{proposition}
\begin{proof}
    We know
    \[
    \phi'(t) = -t\phi(t)
    \]
    and then
    \[
    \overline{\Phi}(t) = \int_t^{\infty}\phi(x)dx \leq \int_t^{\infty} \dfrac{x}{t}\phi(x)dx = \dfrac{1}{t}\phi(x)|_{\infty}^{t} = \dfrac{\phi(t)}{t}, 
    \]
    with
    \[
     \overline{\Phi}(t) = \int_t^{\infty} \phi(x)dx  = \left.\dfrac{\phi(x)}{x}\right|^t_{\infty} - \int_t^{\infty} \dfrac{\phi(x)}{x^2}dx  \geq  \dfrac{\phi(t)}{t} - \dfrac{\overline{\Phi}(t)}{t^2}
    \]
    and we are done.
\end{proof}

\begin{corollary}
    The moment-generating function gives the convenient, slightly looser bound
\[
\mathbb{P}(|G|\ge t)
\le 2e^{-t^2/2},
\qquad t\ge 0. (\text{The inequality is trivial for }t\leq 1)
\]
\end{corollary}

\begin{corollary}
    Let $G_1,\ldots,G_d$ be standard normal random variables. By a union bound,
\[
\mathbb{P}\left(
    \max_{1\le j\le d}|G_j|\ge t
\right)
\le 2d e^{-t^2/2}.
\]

Therefore, with probability at least $1-\delta$,
\[
\max_{1\le j\le d}|G_j|
\le
\sqrt{2\log\left(\frac{2d}{\delta}\right)}.
\]

The scale $\sqrt{\log d}$ is a recurring high-dimensional noise threshold.
\end{corollary}


\begin{definition}[KL Divergence]
Let $P,Q$ be probability measures on $(\mathbb{R}^d,\mathcal{B}^d)$
with $P\ll Q$, and let
\[
r=\frac{\mathrm{d}P}{\mathrm{d}Q}
\]
be their Radon--Nikodym derivative. The Kullback--Leibler (KL) divergence is
\[
D_{\mathrm{KL}}(P\Vert Q)
=
\int \log r\,\mathrm{d}P
=
\int r\log r\,\mathrm{d}Q,
\]
where $0\log 0=0$. For $P\not\ll Q$, the value is defined by $0$.
\end{definition}

\begin{definition}[Mutual information]
The mutual information of $X,Y$ is
\[
I(X;Y)
=
D_{\mathrm{KL}}
\bigl(\mu_{X,Y}\Vert \mu_X\otimes\mu_Y\bigr).
\]
\end{definition}

\begin{definition}[Conditional law and conditional mean]
For random vectors $X\in\mathbb{R}^p$ and $Y\in\mathbb{R}^q$, a
\emph{regular conditional distribution} of $X$ given $Y$ is a family
$K(y,\cdot)$ of probability measures on
$(\mathbb{R}^p,\mathcal{B}(\mathbb{R}^p))$ such that
\[
y\longmapsto K(y,A)
\]
is Borel measurable for every Borel set $A$, and
\begin{equation}
\mathbb{P}(X\in A,Y\in B)
=
\int_B K(y,A)\,\mu_Y(\mathrm{d}y),
\qquad
A\in\mathcal{B}(\mathbb{R}^p),\quad
B\in\mathcal{B}(\mathbb{R}^q).
\end{equation}

We write
\[
K(y,A)=\mathbb{P}(X\in A\mid Y=y).
\]

When $X$ is real and integrable, define
\[
h(y):=\int_{\mathbb{R}}x\,K(y,\mathrm{d}x)
\]
where the absolute integral is finite, and set $h(y)=0$ elsewhere. Then
\[
\mathbb{E}[X\mid Y=y]=h(y)
\qquad\text{and}\qquad
\mathbb{E}[X\mid Y]=h(Y).
\]
\end{definition}

\begin{proposition}
Suppose $X,Y$ have a joint Lebesgue density, or a joint probability
mass function on a countable product support. Conditional laws are
specified for $\mu_X$-almost every $x$, not necessarily at every
conceivable value. Then
\begin{equation}
I(X;Y)
=
\int
D_{\mathrm{KL}}
\bigl(\mu_{Y\mid X=x}\Vert\mu_Y\bigr)
\,\mu_X(\mathrm{d}x).
\end{equation}
\end{proposition}
\begin{proof}
Let
\[
r(x,y)
:=
\frac{\mathrm{d}\mu_{X,Y}}
     {\mathrm{d}(\mu_X\otimes\mu_Y)}(x,y).
\]
It suffices to show that, for $\mu_X$-almost every $x$,
\[
\frac{\mathrm{d}\mu_{Y\mid X=x}}
     {\mathrm{d}\mu_Y}(y)
=
r(x,y)
\qquad
\text{for $\mu_Y$-almost every $y$}.
\]

In the density case, on the set where $p_X(x)p_Y(y)>0$,
\[
\begin{aligned}
\frac{\mathrm{d}\mu_{Y\mid X=x}}
     {\mathrm{d}\mu_Y}(y)
&=
\frac{p_{Y\mid X=x}(y)}{p_Y(y)}\\
&=
\frac{p_{X,Y}(x,y)}
     {p_X(x)p_Y(y)}\\
&=
\frac{\mathrm{d}\mu_{X,Y}}
     {\mathrm{d}(\mu_X\otimes\mu_Y)}(x,y).
\end{aligned}
\]
The same calculation applies in the discrete case, with integrals
replaced by sums.

Consequently,
\begin{align*}
\int
D_{\mathrm{KL}}
\bigl(\mu_{Y\mid X=x}\Vert\mu_Y\bigr)
\,\mu_X(\mathrm{d}x)
&=
\int_{\mathbb{R}^p}
\int_{\mathbb{R}^q}
r(x,y)\log r(x,y)\,
\mu_Y(\mathrm{d}y)\mu_X(\mathrm{d}x)\\
&=
\int_{\mathbb{R}^p\times\mathbb{R}^q}
r\log r\,
\mathrm{d}(\mu_X\otimes\mu_Y)\\
&=
D_{\mathrm{KL}}
\bigl(\mu_{X,Y}\Vert\mu_X\otimes\mu_Y\bigr)\\
&=
I(X;Y).
\end{align*}

If either side is infinite, the argument is justified by applying
Tonelli's theorem to
\[
\psi(r)=r\log r-r+1\geq0,
\]
using
\[
\int r\,\mathrm{d}(\mu_X\otimes\mu_Y)=1.
\]
\end{proof}

\section{Statistical Decision Theory and Sufficiency}

\subsection{Fundamental Definitions}

\begin{definition}
    Let
\[
X:(\Omega,\mathcal{F})\longrightarrow(\mathcal{X},\mathcal{B})
\]
be the random observation. Its distribution is a probability measure on
the sample space $\mathcal{X}$; the realized data are a particular value
$x$. We use $P_\theta$ for the law of the whole observation when
describing a model indexed by an unknown parameter $\theta$.

For an i.i.d.\ sample, write
\[
X_1,\ldots,X_n\sim Q_\theta
\quad\text{independently},
\]
so the full sample has law
\[
P_\theta=Q_\theta^{\otimes n}.
\]
Here $Q_\theta$ is the \textbf{population distribution} of one unit, and
$n$ is the \textbf{sample size}. If each unit has $d$ recorded features,
then
\[
x=(x_1,\ldots,x_n)^\top\in\mathbb{R}^{n\times d}
\]
is a data matrix.

An observation may be a person, a measurement, or a more complicated
object such as an entire graph. The statistical goal is to learn about
the population, or a feature of it, from the observed sample.
\end{definition}

\begin{definition}[Parametric model]
A \emph{parametric model} is a family
\[
\mathcal{P}
=
\left\{
P_\theta:\theta\in\Theta\subseteq\mathbb{R}^k
\right\}
\]
with a fixed finite-dimensional parameter $\theta$. The set $\Theta$ is
the \emph{parameter space}. Inference may concern the whole parameter or
a specified function $\psi(\theta)$ of it.
\end{definition}

\begin{definition}[Dominated model]
A family is \emph{dominated} by a common sigma-finite measure $\nu$ if
\[
P_\theta\ll\nu
\qquad\text{for every }\theta.
\]
Then its densities are
\[
p_\theta=\frac{\mathrm{d}P_\theta}{\mathrm{d}\nu}.
\]
Lebesgue measure gives probability density functions; counting measure
gives probability mass functions. The reference measure is fixed across
the family.
\end{definition}

\begin{definition}[Identifiability]
A parametrization
\[
\{P_\theta:\theta\in\Theta\}
\]
is \emph{identifiable} if
\[
P_{\theta_1}=P_{\theta_2}
\quad\Longrightarrow\quad
\theta_1=\theta_2.
\]

Even in a non-identifiable parametrization, a target $\psi(\theta)$ can
be identifiable: it must have the same value whenever two parameters
give the same observation law.
\end{definition}

\begin{definition}[Statistic]
A \emph{statistic} is a measurable function $T(X)$ of the observed data.
The rule
\[
T:\mathcal{X}\longrightarrow\mathcal{T}
\]
must be specified without using an unknown parameter. Its distribution
may nevertheless depend on that parameter. For real-vector data and
summaries, measurability is with respect to their Borel
sigma-algebras.
\end{definition}

\begin{definition}
    [Order Statistic]The order statistics
\[
X_{(1)}\leq\cdots\leq X_{(n)}
\]
are the sample values sorted in increasing order, retaining repeats if
there are ties. Suppose the observations are i.i.d.\ with CDF $F$ and
Lebesgue density $f$. For $1\leq i\leq n$, the $i$th order statistic has
density
\begin{equation}
g_i(x)
=
\frac{n!}{(i-1)!(n-i)!}
F(x)^{i-1}\bigl\{1-F(x)\bigr\}^{n-i}f(x),
\qquad\text{for Lebesgue-a.e.\ }x.
\end{equation}
\end{definition}

\begin{definition}[Exponential family]
A dominated family is an \emph{exponential family} if its densities have
the form
\begin{equation}
p_\theta(x)
=
\exp\left\{
\eta(\theta)^\top T(x)-\xi(\theta)
\right\}h(x),
\qquad
\xi(\theta)=A\bigl(\eta(\theta)\bigr),
\tag{5}
\end{equation}
where $T(x)\in\mathbb{R}^m$ and $h(x)\geq0$ do not depend on $\theta$,
and
\[
A(\eta)
=
\log\int
\exp\left\{\eta^\top T(x)\right\}
h(x)\,\mathrm{d}\nu(x).
\]

The \emph{natural parameter} is $\eta$. Its natural parameter space
$\mathcal{H}$ consists of those $\eta$ for which this integral is
positive and finite. In canonical form,
\[
p_\eta(x)
=
\exp\left\{
\eta^\top T(x)-A(\eta)
\right\}h(x).
\]
The function $A$ is the \emph{log-normalizing function}.
\end{definition}